% !PRINTED 539
\weberpage{539}

%<<CONT
\weq{16}{\begin{aligned}
\operatorname*{Lim}_{s=1} (4^s + 2 - 3.\,2^s) \overset{\psi}{\Sigma} S &= \frac{24\pi}{\sqrt{m}} \log \Pi \frac{f(\omega)}{\sqrt[4]{2}}, & m &\equiv 3 \pmod 8, \\
\operatorname*{Lim}_{s=1} (4^s + 2 - 3.\,2^s) \overset{\psi}{\Sigma} S &= \frac{8\pi}{\sqrt{m}} \log \Pi \frac{f(\omega)}{\sqrt[4]{2}}, & m &\equiv 7 \pmod 8.
\end{aligned}}
Es ist aber [§ 141, (28)]:
\[
\mathrm{Lim}\,(s-1) \overset{\psi}{\Sigma} S = \frac{\pi h}{\sqrt{m}}, \quad \mathrm{Lim}\,\frac{4^s + 2 - 3.\,2^s}{s-1} = 2 \log 2,
\]
also:
\weq{17}{\begin{aligned}
\log \Pi \frac{f(\omega)}{\sqrt[3]{2}} &= 0, & m &\equiv 3 \pmod 8, \\
\log \Pi \frac{f(\omega)}{\sqrt{2}} &= 0, & m &\equiv 7 \pmod 8.
\end{aligned}}
Die erste dieser Formeln gibt das bereits auf andere Weise
abgeleitete Resultat; die zweite gibt das neue:
\weq{18}{\tau = \tfrac{1}{2}, \quad m \equiv 7 \pmod 8.}
4. Es bleibt noch die Bestimmung von $\tau$ in dem Falle übrig,
wo $m$ durch eine höhere Potenz von 2 teilbar ist. Um auch noch
diese Bestimmung auszuführen, sei
\weq{19}{m = 4m',}
% [sic] The source prints sqrt(m); (19), m'=AC-B^2 and the root of (A,2B,C) require sqrt(m') here.
\weq{20}{\omega' = \frac{-B + i\sqrt{m}}{A}, \quad m' = AC - B^2,}
also $\omega'$ die Wurzel der quadratischen Form $(A, 2B, C)$ der
Diskriminante $-m$, worin $A$ und $C$ als \gesperrt{ungerade} vorausgesetzt
werden können. Es sind dann
\weq{21}{\omega_1 = 2\,\omega', \quad \omega_2 = \frac{\omega'}{2}}
Wurzeln der Formen
\weq{22}{(A, 4B, 4C), \quad (4A, 4B, C),}
und es sind nach § 126, (13) die 24sten Potenzen von
\[
f_1(\omega_1) = \frac{\sqrt{2}}{f_2(\omega')}, \quad f_2(\omega_2) = \frac{\sqrt{2}}{f_1(\omega')}
\]
Klasseninvarianten der Diskriminante $-4m$.

Durchläuft $\omega'$ ein vollständiges Wurzelsystem der Diskriminante $-m$, so durchlaufen $\omega_1$ und $\omega_2$ zusammen ein vollständiges
Wurzelsystem der Diskriminante $-4m$ (§ 123), und wir erhalten,
%CONT>>
